% ═══════════════════════════════════════════════════════════════
%  Abstract (English & French)
% ═══════════════════════════════════════════════════════════════
\chapter*{Abstract}
\addcontentsline{toc}{chapter}{Abstract}

\noindent\textbf{Keywords:} Snowflake, CI/CD, DataOps, DevOps, Environment Synchronization, Secure Data Sharing, GitHub Actions, Data Engineering, Data Governance, IBM, ARYZTA.

\vspace{0.5cm}

This report presents the technical and methodological work carried out during a final-year internship at the \textbf{IBM Client Innovation Center (CIC) Morocco},a Snowflake-focused data engineering company recently acquired by IBM. The project was conducted for \textbf{ARYZTA}, one of the world's largest industrial bakery companies, whose analytical and operational reporting platform is built on \textbf{Snowflake}.

The primary challenge addressed in this project was the critical misalignment between ARYZTA's three Snowflake environments: Development (\texttt{DEV}), User Acceptance Testing (\texttt{UAT}), and Production (\texttt{PROD}). Due to the absence of DevOps practices and version control, significant discrepancies had accumulated over several years, including missing database objects, inconsistent schemas, and divergent pipeline configurations. As the lower environments no longer accurately reflected production, development teams were forced to implement changes directly in the production environment, increasing both platform drift and the risk of operational incidents.

To address these challenges, the project was structured around three complementary contributions. First, a \textbf{read-only auditing tool} was designed and developed to compare two Snowflake environments and accurately identify and quantify their differences without modifying either environment. Second, a \textbf{multi-environment synchronization strategy} based on Snowflake's \textbf{Secure Data Sharing} mechanism was implemented to synchronize the \texttt{UAT} and \texttt{DEV} environments with \texttt{PROD}, which serves as the single source of truth. Finally, a \textbf{CI/CD pipeline} integrated with \textbf{GitHub} and automated through \textbf{GitHub Actions} was established to govern deployments using reviewed pull requests, automated validation stages, and environment-specific promotion gates, thereby eliminating manual deployments to production.

The project was carried out following an \textbf{Agile} methodology, with task management supported by \textbf{ServiceNow} and collaborative development managed through \textbf{GitHub}. The resulting solution significantly improves the governance, reliability, and maintainability of ARYZTA's Snowflake data platform while establishing a robust DevOps foundation for future development and deployment activities.

\vspace{1cm}
\hrule
\vspace{0.5cm}


\chapter*{Résumé}
\addcontentsline{toc}{chapter}{Résumé}

\noindent\textbf{Mots-clés:} Snowflake, CI/CD, DataOps, DevOps, Synchronisation d'environnements, Partage Sécurisé de Données, GitHub Actions, Ingénierie des données, Gouvernance, IBM, ARYZTA.

\vspace{0.5cm}

Ce rapport documente les travaux techniques et méthodologiques réalisés dans le cadre d'un stage de fin d'études effectué au sein de l'\textbf{IBM Client Innovation Center (CIC) Maroc}, en collaboration avec \textbf{Hakkoda}, une société spécialisée dans l'ingénierie des données sur Snowflake récemment acquise par IBM. Le projet a été mené pour le compte d'\textbf{ARYZTA}, l'une des plus grandes entreprises de boulangerie industrielle au monde, dont la plateforme de données repose sur \textbf{Snowflake} pour ses activités analytiques et de reporting opérationnel.

Le principal défi traité concerne l'état critique de désalignement entre les trois environnements Snowflake d'ARYZTA --- Développement (\texttt{DEV}), Test d'Acceptation Utilisateur (\texttt{UAT}) et Production (\texttt{PROD}). En l'absence de pratiques DevOps et de tout contrôle de version, des divergences majeures se sont accumulées au fil des années : objets de base de données manquants, schémas incohérents et configurations de pipelines divergentes. Faute d'environnements inférieurs représentatifs, les équipes étaient contraintes d'appliquer leurs modifications directement en production, accentuant la dérive et le risque d'incidents opérationnels.

Pour répondre à cette problématique, le travail s'articule autour de trois contributions complémentaires. Premièrement, un \textbf{outil d'audit en lecture seule} a été conçu et développé afin de comparer deux environnements Snowflake et de quantifier précisément leurs écarts, sans jamais pouvoir les modifier. Deuxièmement, une \textbf{stratégie de synchronisation multi-environnements} fondée sur le \textbf{Partage Sécurisé de Données (Secure Data Sharing)} de Snowflake réaligne quotidiennement les environnements \texttt{UAT} et \texttt{DEV} sur \texttt{PROD}, considérée comme la source de référence. Troisièmement, un \textbf{pipeline CI/CD} intégré à \textbf{GitHub} et exécuté via \textbf{GitHub Actions} encadre les déploiements par des demandes de tirage (\textit{pull requests}) revues, des étapes de validation automatisées et des portes de promotion propres à chaque environnement, supprimant ainsi tout déploiement manuel en production.

Le projet a été conduit selon une démarche \textbf{Agile}, avec un suivi des tâches assuré par \textbf{ServiceNow} et un développement collaboratif géré sous \textbf{GitHub}. Les résultats obtenus constituent une amélioration significative de la gouvernance de la plateforme de données d'ARYZTA, ainsi que de la fiabilité et de la maintenabilité de ses déploiements.

\newpage
